\documentclass[a4paper,10pt]{report}\input{../0.Préambule/Préambule_cours_prof}\begin{document}

\newpage
\pagestyle{empty}
\section*{Activité Géogébra : Les parallélogrammes}

\paragraph{Construction 1}:

\begin{prop}
Si un quadrilatère a ses côtés opposés parallèles, alors c'est un parallélogramme. 
\end{prop}

\medskip
\noindent Ouvrir le fichier Géogébra "Construction 1.ggb" donné par le professeur.

\medskip
\noindent En utilisant la propriété 1, construire, à partir des points $A$, $B$ et $C$ déjà placés, le parallélogramme $ABCD$.

\medskip
\noindent Enregistrer le travail.

\paragraph{Construction 2}:

\medskip
\noindent Ouvrir le fichier Géogébra "Construction 2.ggb" donné par le professeur.

\medskip
\noindent En utilisant la définition 1, construire, à partir des points $A$, $B$ et $C$ déjà placés, le parallélogramme $ABCD$.

\medskip
\noindent Enregistrer le travail.

\paragraph{Construction 3}:

\begin{prop}
Si un quadrilatère non croisé a des côtés opposés de même longueur, alors c'est un parallélogramme. 
\end{prop}

\medskip
\noindent Ouvrir le fichier Géogébra "Construction 3.ggb" donné par le professeur.

\medskip
\noindent En utilisant la propriété 2, construire, à partir des points $A$, $B$ et $C$ déjà placés, le parallélogramme $ABCD$.

\medskip
\noindent Enregistrer le travail.

\paragraph{Construction 4}:

\begin{prop}
Si un quadrilatère a deux côtés opposés parallèles et de même longueur, alors c'est un parallélogramme. 
\end{prop}

\medskip
\noindent Ouvrir le fichier Géogébra "Construction 4.ggb" donné par le professeur.

\medskip
\noindent En utilisant la propriété 3, construire, à partir des points $A$, $B$ et $C$ déjà placés, le parallélogramme $ABCD$.

\medskip
\noindent Enregistrer le travail.

\paragraph{Construction 5}:

\begin{prop}
Si un quadrilatère a ses diagonales qui se coupent en leur milieu, alors c'est un parallélogramme. 
\end{prop}

\medskip
\noindent Ouvrir le fichier Géogébra "Construction 5.ggb" donné par le professeur.

\medskip
\noindent En utilisant la propriété 4, construire, à partir des points $A$, $B$ et $C$ déjà placés, le parallélogramme $ABCD$.

\medskip
\noindent Enregistrer le travail.

\paragraph{Construction 6}:

\medskip
\noindent Ouvrir le fichier Géogébra "Construction 6.ggb" donné par le professeur.

\medskip
\noindent En utilisant la propriété 4, construire, à partir des points $A$, $B$ et $C$ déjà placés, le parallélogramme $ABCD$.

\medskip
\noindent Enregistrer le travail.

\end{document}